% Vietnamese translation for OI Public Library
% Source-PDF: IOI中国国家候选队论文/2009/毛杰明/母函数的性质及应用.pdf
\documentclass[11pt]{article}
\usepackage{oipl-vn}

\newcommand{\Oh}{\mathcal{O}}
\newcommand{\seq}[1]{\left\langle #1\right\rangle}
\newcommand{\coef}[2]{[x^{#2}]\,#1}
\newcommand{\egfcoef}[2]{[x^{#2}/#2!]\,#1}
\newcommand{\C}[2]{\binom{#1}{#2}}
\newcommand{\eps}{\varepsilon}

\title{Tính chất và ứng dụng của hàm sinh}
\author{Mao Jieming\\Trường Ngoại ngữ Nam Kinh\\Email: \texttt{maojm517@163.com}\\Mã bưu chính: 210008}
\date{Bài luận đội tuyển quốc gia năm 2009}

\begin{document}
\maketitle

\origtitle{Tính chất và ứng dụng của hàm sinh}
\sourcepath{IOI China candidate team papers / 2009 / Mao Jieming / generating functions}

\begin{abstract}
Lý thuyết hàm sinh là một phương pháp quan trọng của toán rời rạc, là cầu nối
giữa toán rời rạc và toán liên tục. Nó là công cụ mạnh để xử lý dãy số và các
bài toán đếm tổ hợp. Sức mạnh của nó nằm ở chỗ có thể dùng một cách thống
nhất để giải nhiều loại bài toán khác nhau. Vì vậy, khi giải một số bài trong
thi tin học, hàm sinh có hiệu quả đáng ngạc nhiên; nó cũng có vai trò quan
trọng trong tối ưu hóa thuật toán và chứng minh mệnh đề. Bài viết trước hết
giới thiệu các tính chất của hàm sinh, sau đó dùng một số bài toán điển hình
để trình bày ứng dụng của hàm sinh trong thi tin học, cuối cùng đưa ra một
tổng kết về hàm sinh.
\end{abstract}

\noindent\textbf{Từ khóa:} hàm sinh, truy hồi, hoán vị và tổ hợp.

\tableofcontents

\section{Tính chất của hàm sinh}

\subsection{Định nghĩa}

Hàm sinh là chuỗi lũy thừa dùng để tương ứng với một dãy vô hạn. Nói chung,
hàm sinh có dạng
\[
  G(x)=g_0+g_1x+g_2x^2+\cdots=\sum_{n\ge 0} g_n x^n.
\]
Ta gọi \(G(x)\) là hàm sinh của dãy
\(\seq{g_0,g_1,g_2,\ldots}\), và ký hiệu
\[
  G(x)\longleftrightarrow \seq{g_0,g_1,g_2,\ldots}.
\]

Ví dụ, dãy \(\seq{1,1,1,1,1,\ldots}\) có hàm sinh
\[
  G(x)=1+x+x^2+x^3+\cdots=\sum_{n\ge 0}x^n.
\]
Theo công thức tổng cấp số nhân,
\[
  G(x)=1+x+x^2+x^3+\cdots=\frac{1}{1-x}.
\]
Ở đây ta gọi \(\frac1{1-x}\) là dạng đóng của hàm sinh này.

Cần chú ý rằng do vấn đề hội tụ, công thức trên không đúng khi \(x\ge 1\).
Nhưng khi xét hàm sinh, ta không cần thay giá trị cụ thể của \(x\) để tính
giá trị hàm, nên thông thường không cần xét tính hội tụ. Bài viết này cũng bỏ
qua vấn đề hội tụ của hàm sinh.

\subsection{Các thao tác cơ bản}

\paragraph{1. Nhân với hằng số.}
\[
  cG(x)=cg_0+cg_1x+cg_2x^2+\cdots
  \longleftrightarrow \seq{cg_0,cg_1,cg_2,\ldots}.
\]

\paragraph{2. Cộng và trừ.}
\[
\begin{aligned}
  F(x)\pm G(x)
    &=(f_0\pm g_0)+(f_1\pm g_1)x+(f_2\pm g_2)x^2+\cdots\\
    &\longleftrightarrow
      \seq{f_0\pm g_0,f_1\pm g_1,f_2\pm g_2,\ldots}.
\end{aligned}
\]

\paragraph{3. Dịch phải.}
Thao tác này dịch dãy sang phải \(k\) vị trí và bù \(k\) số \(0\) ở đầu:
\[
  x^kG(x)=g_0x^k+g_1x^{k+1}+g_2x^{k+2}+\cdots
  \longleftrightarrow
  \seq{\underbrace{0,0,\ldots,0}_{k\text{ số }0},g_0,g_1,g_2,\ldots}.
\]

\paragraph{4. Lấy đạo hàm.}
\[
  G'(x)=g_1+2g_2x+3g_3x^2+\cdots
  \longleftrightarrow \seq{g_1,2g_2,3g_3,\ldots}.
\]
Lấy đạo hàm có hai hiệu ứng: dịch dãy sang trái một vị trí, đồng thời nhân
mỗi hạng với chỉ số ban đầu của nó.

Ví dụ, đặt
\[
  G(x)=1+x+x^2+x^3+\cdots=\frac1{1-x}.
\]
Khi đó
\[
  \frac1{(1-x)^2}=G'(x)=1+2x+3x^2+4x^3+\cdots.
\]
Vì vậy dạng đóng \(\frac1{(1-x)^2}\) tương ứng với dãy
\(\seq{1,2,3,4,5,\ldots}\).

\paragraph{5. Quy tắc tích chập.}
\[
\begin{aligned}
  H(x)&=G(x)\cdot F(x)\\
      &=(g_0+g_1x+g_2x^2+\cdots)(f_0+f_1x+f_2x^2+\cdots).
\end{aligned}
\]
Khi đó hệ số của \(x^n\) trong \(H(x)\) là
\[
  h_n=g_0f_n+g_1f_{n-1}+g_2f_{n-2}+\cdots+g_nf_0.
\]
Từ đây cũng có thể chứng minh rằng dạng đóng \(\frac1{(1-x)^2}\) tương ứng
với dãy \(\seq{1,2,3,4,5,\ldots}\).

Điểm quan trọng của quy tắc tích chập trong bài toán tổ hợp là: nếu \(G(x)\)
tương ứng với hàm sinh của cách chọn phần tử trong \(A\), còn \(F(x)\) tương
ứng với cách chọn phần tử trong \(B\), thì \(G(x)F(x)\) tương ứng với cách
chọn phần tử trong \(A\cup B\).

\subsection{Hàm sinh của một số dãy đơn giản}

Từ các thao tác cơ bản ở trên, ta đã có thể tìm hàm sinh của nhiều dãy.
Phía trên ta đã xét các dạng đóng \(\frac1{1-x}\) và
\(\frac1{(1-x)^2}\). Bây giờ xét một hàm sinh thường dùng:
\[
  \frac1{(1-x)^m}.
\]
Ta có
\[
  G(x)=\frac1{(1-x)^m}
      =\underbrace{\frac1{1-x}\cdots\frac1{1-x}}_{m\text{ lần}}
      =(1+x+x^2+\cdots)^m.
\]
Do đó hệ số \(g_n\) của \(x^n\) trong \(G(x)\) bằng số nghiệm nguyên không âm
của phương trình bất định
\[
  x_1+x_2+\cdots+x_m=n.
\]
Đây là bài toán quen thuộc; dùng phương pháp ``thanh chắn và quả bóng'' ta có
\[
  g_n=\C{m+n-1}{m-1}.
\]
Cũng có thể chứng minh bằng quy nạp từ
\[
  \frac1{(1-x)^m}=\frac1{1-x}\cdot\frac1{(1-x)^{m-1}}.
\]
Vì vậy
\[
  \frac1{(1-x)^m}
  \longleftrightarrow
  \seq{1,\C{m}{m-1},\C{m+1}{m-1},\C{m+2}{m-1},\ldots}.
\]

Tiếp theo xét dạng đóng \(\frac1{1-\alpha x}\). Có thể xem \(\alpha x\) như
một biến duy nhất rồi khai triển:
\[
  \frac1{1-\alpha x}
  =1+\alpha x+\alpha^2x^2+\alpha^3x^3+\cdots
  \longleftrightarrow \seq{1,\alpha,\alpha^2,\ldots}.
\]
Kết hợp với kết quả của \(\frac1{(1-x)^m}\), ta được kết luận hữu dụng:
\[
  \frac1{(1-\alpha x)^m}
  =1+\alpha\C{m}{m-1}x+\alpha^2\C{m+1}{m-1}x^2
   +\alpha^3\C{m+2}{m-1}x^3+\cdots.
\]

\subsection{Hàm sinh mũ}

Đôi khi hàm sinh thường của dãy \(\seq{g_n}\) có tính chất rất phức tạp,
nhưng hàm sinh của dãy \(\seq{g_n/n!}\) lại rất đơn giản. Khi đó ta muốn
nghiên cứu \(\seq{g_n/n!}\) trước, rồi nhân lại với \(n!\).

Ta gọi
\[
  G(x)=\sum_{n\ge 0} g_n\frac{x^n}{n!}
\]
là hàm sinh mũ của dãy \(\seq{g_0,g_1,g_2,\ldots}\). Vì hàm sinh mũ thỏa quy
tắc tích chập phù hợp với bài toán hoán vị, nên ta thường dùng hàm sinh mũ để
xử lý các bài đếm có xét thứ tự.

Trước hết xét dãy cơ bản \(\seq{1,1,1,1,\ldots}\). Theo định nghĩa, hàm sinh
mũ của nó là
\[
  G(x)=1+\frac{x}{1!}+\frac{x^2}{2!}+\frac{x^3}{3!}+\frac{x^4}{4!}+\cdots.
\]
Dạng đóng của \(G(x)\) không dễ thu được chỉ bằng các phương pháp sơ cấp; ở
đây dùng công thức Taylor:
\[
  f(x)=f(x_0)+f'(x_0)(x-x_0)+\frac{f''(x_0)}{2!}(x-x_0)^2
       +\frac{f'''(x_0)}{3!}(x-x_0)^3+\cdots.
\]
Lấy \(f(x)=e^x\), \(x_0=0\), ta được
\[
  e^x=1+x+\frac{x^2}{2!}+\frac{x^3}{3!}+\frac{x^4}{4!}+\cdots=G(x).
\]
Nói cách khác, dạng đóng của hàm sinh mũ của dãy
\(\seq{1,1,1,1,\ldots}\) là \(e^x\).

Tương tự, dùng công thức Taylor ta có:
\[
  \seq{1,-1,1,-1,1,-1,\ldots}
  \quad\text{có hàm sinh mũ}\quad e^{-x},
\]
\[
  \seq{0,1,0,1,0,1,\ldots}
  \quad\text{có hàm sinh mũ}\quad \frac{e^x-e^{-x}}2,
\]
\[
  \seq{1,0,1,0,1,0,\ldots}
  \quad\text{có hàm sinh mũ}\quad \frac{e^x+e^{-x}}2.
\]

Các thao tác nhân với hằng số, cộng và trừ của hàm sinh mũ giống với hàm sinh
thường, nhưng dịch trái và dịch phải thì hơi khác. Giả sử
\[
  G(x)=g_0+\frac{g_1}{1!}x+\frac{g_2}{2!}x^2+\cdots.
\]
Khi đó
\[
  G'(x)=g_1+\frac{g_2}{1!}x+\frac{g_3}{2!}x^2+\cdots.
\]
Vì vậy đạo hàm thực hiện thao tác dịch trái đối với hàm sinh mũ. Tương tự,
tích phân thực hiện thao tác dịch phải.

Giống hàm sinh thường, phép toán thú vị nhất của hàm sinh mũ là phép nhân.
Giả sử \(G(x)\) là hàm sinh mũ của dãy \(\seq{g_0,g_1,g_2,\ldots}\), còn
\(F(x)\) là hàm sinh mũ của dãy \(\seq{f_0,f_1,f_2,\ldots}\):
\[
  G(x)=g_0+\frac{g_1}{1!}x+\frac{g_2}{2!}x^2+\cdots,\qquad
  F(x)=f_0+\frac{f_1}{1!}x+\frac{f_2}{2!}x^2+\cdots.
\]
Nếu \(H(x)=G(x)F(x)\), thì
\[
\begin{aligned}
  H(x)
  &=\sum_{n\ge0}\frac{x^n}{n!}
    \sum_{k=0}^n
      \frac{g_k f_{n-k} n!}{k!(n-k)!}\\
  &=\sum_{n\ge0}\frac{x^n}{n!}
    \sum_{k=0}^n g_k f_{n-k}\C{n}{k}.
\end{aligned}
\]
Vì vậy
\[
  h_n=\sum_{k=0}^n g_k f_{n-k}\C{n}{k}.
\]
Khi xét việc lấy \(k\) phần tử từ một loại và \(n-k\) phần tử từ loại còn
lại, có \(\C{n}{k}\) cách sắp vị trí cho hai loại phần tử đó. Đây chính là lý
do hàm sinh mũ thích hợp cho bài toán hoán vị.

\subsection{Định lý Pólya dạng hàm sinh}

Định lý Pólya hiện đã là một định lý quen thuộc. Các bài viết của Fu Wenjie
năm 2001 và Chen Yuxi năm 2008 đều giới thiệu định lý này. Phát biểu cơ bản
như sau.

\term{Định lý Pólya.}
Giả sử \(G\) là một nhóm hoán vị trên \(n\) đối tượng. Nếu dùng \(m\) màu để
tô \(n\) đối tượng này, thì số phương án tô màu khác nhau là
\[
  l=\frac1{|G|}\sum_{a\in G}m^{c(a)},
\]
trong đó \(c(a)\) là số chu trình của hoán vị \(a\).

Định lý này chỉ dùng để đếm số lượng, chứ không cho biết chi tiết cấu trúc
của các trạng thái.

Ví dụ, dùng hai màu \(b,w\) để tô \(4\) viên bi. Nếu phân biệt từng viên bi
thì có \(16\) trường hợp:
\[
\begin{gathered}
bbbb,\ bbbw,\ bbwb,\ bbww,\ bwbb,\ bwbw,\ bwwb,\ bwww,\\
wbbb,\ wbbw,\ wbwb,\ wbww,\ wwbb,\ wwbw,\ wwwb,\ wwww.
\end{gathered}
\]
Nếu không quan tâm viên bi cụ thể nào nhận màu nào, chẳng hạn \(wbbw\) được
viết gọn thành \(b^2w^2\), thì toàn bộ phương án có thể viết thành
\((b+w)^4\), vì \(4\) viên bi không phân biệt và phép nhân giao hoán:
\[
  (b+w)^4=b^4+4b^3w+6b^2w^2+4bw^3+w^4.
\]
Tư tưởng này cho phép mở rộng định lý Pólya, gọi là định lý Pólya dạng hàm
sinh.

Trong công thức Pólya, thay \(m^{c(a_i)}\) bằng
\[
  (b_1+b_2+\cdots+b_m)^{c_1(a_i)}
  (b_1^2+b_2^2+\cdots+b_m^2)^{c_2(a_i)}
  \cdots
  (b_1^n+b_2^n+\cdots+b_m^n)^{c_n(a_i)},
\]
ta được
\[
  P=\frac1{|G|}\sum_{a_i\in G}
  \prod_{j=1}^n
  (b_1^j+b_2^j+\cdots+b_m^j)^{c_j(a_i)}.
\]
Ở đây \(c_j(a_i)\) là số chu trình có độ dài \(j\) trong hoán vị \(a_i\).

Ví dụ, có \(4\) hạt \(v_1,v_2,v_3,v_4\) xếp thành vòng tròn, dùng hai màu
\(b,w\) để tô; các phương án trùng nhau sau phép quay được xem là một. Nhóm
hoán vị \(G\) gồm
\[
  (v_1)(v_2)(v_3)(v_4),\quad
  (v_1v_2v_3v_4),\quad
  (v_1v_3)(v_2v_4),\quad
  (v_1v_4v_3v_2).
\]
Dùng định lý Pólya thường, số phương án khác nhau là
\[
  l=\frac14(2^4+2+2^2+2)=6.
\]
Dùng định lý Pólya dạng hàm sinh, ta thu được các phương án cụ thể:
\[
\begin{aligned}
  P
  &=\frac14\left[(b+w)^4+(b^4+w^4)+(b^2+w^2)^2+(b^4+w^4)\right]\\
  &=b^4+b^3w+2b^2w^2+bw^3+w^4.
\end{aligned}
\]
Nghĩa là trường hợp \(2\) đen \(2\) trắng có hai loại; các trường hợp
\(4\) đen, \(4\) trắng, \(3\) đen \(1\) trắng và \(3\) trắng \(1\) đen đều
chỉ có một loại.

\section{Ứng dụng của hàm sinh}

\subsection{Một bài tự đặt}

\term{Ví dụ 1.}
Xiaoming đi du lịch và cần mang một số thức ăn, gồm khoai tây chiên, sô-cô-la,
nước khoáng, hamburger, sữa và kẹo. Sau khi ước lượng, cậu thấy mang
\(n\) món thức ăn là phù hợp, với \(n\le 10^{100}\). Nhưng cậu có vài sở
thích:
\begin{itemize}
  \item mang nhiều nhất \(1\) hamburger;
  \item số miếng sô-cô-la là bội của \(5\);
  \item mang nhiều nhất \(4\) chai nước khoáng;
  \item số gói khoai tây chiên là số chẵn;
  \item mang nhiều nhất \(3\) hộp sữa;
  \item số viên kẹo là bội của \(4\).
\end{itemize}
Hỏi Xiaoming có bao nhiêu cách chuẩn bị thức ăn cho chuyến đi này?

\paragraph{Phân tích.}
Cách nghĩ thông thường là trước hết xét ba điều kiện \(1,3,5\), có thể giải
bằng liệt kê; sau đó bài toán trở thành đếm số nghiệm nguyên không âm của
phương trình
\[
  5x+2y+4z=m.
\]
Cách đầu tiên là đặt \(f(m)\) là số nghiệm nguyên không âm của phương trình
trên, rồi tách ba biến ra để làm ba lần truy hồi \(\Oh(n)\).

Cách khác là liệt kê giá trị của \(x\), sau đó dùng thuật toán Euclid mở rộng
để tìm một nghiệm của phương trình còn lại và tính số nghiệm.

Hai thuật toán này đều không đáp ứng được kích thước \(n\) của đề bài. Cách
thứ hai cũng không có tính mở rộng tốt khi số điều kiện tăng lên.

Ta xem hàm sinh giải bài này như thế nào. Hamburger có thể không mang hoặc
mang một cái, nên hàm sinh là
\[
  h(x)=1+x.
\]
Sô-cô-la phải được mang theo số lượng là bội của \(5\), tức có thể mang
\(0,5,10,\ldots\) miếng:
\[
  c(x)=1+x^5+x^{10}+x^{15}+\cdots=\frac1{1-x^5}.
\]
Khoai tây chiên phải mang theo số gói là bội của \(2\):
\[
  p(x)=1+x^2+x^4+x^6+\cdots=\frac1{1-x^2}.
\]
Nước khoáng mang nhiều nhất \(4\) chai:
\[
  w(x)=1+x+x^2+x^3+x^4=\frac{1-x^5}{1-x}.
\]
Sữa mang nhiều nhất \(3\) hộp:
\[
  m(x)=1+x+x^2+x^3=\frac{1-x^4}{1-x}.
\]
Kẹo được mang theo số lượng là bội của \(4\):
\[
  s(x)=1+x^4+x^8+x^{12}+\cdots=\frac1{1-x^4}.
\]
Theo quy tắc tích chập, hàm sinh của số cách chọn \(n\) món thức ăn là
\[
\begin{aligned}
  h(x)c(x)p(x)w(x)m(x)s(x)
  &=(1+x)\frac1{1-x^5}\frac1{1-x^2}
    \frac{1-x^5}{1-x}\frac{1-x^4}{1-x}\frac1{1-x^4}\\
  &=\frac1{(1-x)^3}\\
  &=1+\C{3}{2}x+\C{4}{2}x^2+\C{5}{2}x^3+\cdots.
\end{aligned}
\]
Hệ số của \(x^n\) là \(\C{n+2}{2}\), chính là đáp án.

Như vậy, dùng hàm sinh, bài toán được giải trong \(\Oh(1)\) và gần như không
có độ khó lập trình; khi \(n\) rất lớn chỉ cần dùng nhân số lớn.

Đáp án của bài này khá đặc biệt, nên có thể có bạn thấy tìm quy luật còn
nhanh hơn. Nhưng với các bài mà kết quả không gọn như vậy, cách tìm quy luật
thường mất tác dụng.

\paragraph{Tiểu kết.}
Bài toán này thể hiện rõ ưu thế của hàm sinh: các điều kiện hạn chế được
chuyển thành những đối tượng dễ xử lý, rồi được giải bằng thao tác rất cơ
học. Qua bài này ta thấy hàm sinh không chỉ có thể giảm độ phức tạp thời gian
và lượng mã, mà còn giảm độ khó tư duy.

\subsection{Chocolate}

\term{Ví dụ 2.}
ACM/ICPC Regional Contest Beijing 2002, Problem F: \textit{Chocolate}.

Trong một túi có sô-cô-la với \(c\) màu. Mỗi lần lấy một viên từ túi. Nếu viên
vừa lấy có cùng màu với một viên đang ở trên bàn, thì lấy cả hai viên đó đi;
ngược lại đặt viên vừa lấy lên bàn. Giả sử xác suất lấy được mỗi màu là như
nhau. Hãy tính xác suất để sau khi lấy \(n\) viên, trên bàn còn \(m\) viên
sô-cô-la. Giới hạn \(c\le 100\), \(n,m\le 1000000\).

\paragraph{Phân tích.}
Trước hết, chỉ khi \(2\mid (m-n)\), \(m\le n\) và \(m\le c\), đáp án mới khác
\(0\). Dưới đây chỉ xét các trường hợp có thể khác \(0\).

Cách dễ nghĩ là truy hồi. Đặt \(f[i,j]\) là xác suất sau khi lấy \(i\) viên,
trên bàn còn \(j\) viên. Khi đó
\[
  f[i,j]=\frac{c-j+1}{c}f[i-1,j-1]+\frac{j+1}{c}f[i-1,j+1].
\]
Cách này có độ phức tạp \(\Oh(nc)\), rất dễ quá thời gian với giới hạn đề.

Bây giờ xét cách dùng hàm sinh. Ta tính xác suất bằng cách chia mọi tình huống
thành các tình huống đồng xác suất, rồi xem tình huống cần đếm chiếm bao
nhiêu. Vì mỗi lần lấy ra một viên có màu đều nhau, tổng số tình huống theo
dãy \(n\) lần lấy là \(c^n\). Cách chia này có xét thứ tự giữa các màu trong
dãy lấy ra, nên bài toán thích hợp với hàm sinh mũ.

Điều kiện còn \(m\) viên trên bàn tương đương với: trong \(c\) màu, có \(m\)
màu xuất hiện số lần lẻ, và \(c-m\) màu xuất hiện số lần chẵn.

Xét từng màu sô-cô-la. Với các màu xuất hiện số lần lẻ, hàm sinh mũ tương ứng
có hệ số bậc lẻ bằng \(1\) và hệ số bậc chẵn bằng \(0\), tức là hàm sinh mũ
của dãy \(\seq{0,1,0,1,0,1,\ldots}\):
\[
  \frac{e^x-e^{-x}}2.
\]
Tương tự, với các màu xuất hiện số lần chẵn, hàm sinh mũ là
\[
  \frac{e^x+e^{-x}}2.
\]
Vì vậy hàm sinh mũ tương ứng với \(m\) màu xuất hiện lẻ và \(c-m\) màu xuất
hiện chẵn là
\[
  G(x)=\left(\frac{e^x-e^{-x}}2\right)^m
       \left(\frac{e^x+e^{-x}}2\right)^{c-m}.
\]
Nếu \(\egfcoef{G(x)}{n}=g_n\), thì số tình huống cần đếm là
\[
  g_n\cdot n!\cdot \C{c}{m},
\]
trong đó \(n!\) đến từ hàm sinh mũ, còn \(\C{c}{m}\) là số cách chọn \(m\)
màu xuất hiện lẻ trong \(c\) màu. Xác suất cần tìm là
\[
  \frac{g_n n!\C{c}{m}}{c^n}.
\]

Ta còn phải khai triển
\[
  G(x)=\left(\frac{e^x-e^{-x}}2\right)^m
       \left(\frac{e^x+e^{-x}}2\right)^{c-m}.
\]
Xem \(e^x\) như một biến, khai triển \(G(x)\) thành đa thức theo \(e^x\). Với
mỗi hạng \(e^{kx}\), hệ số của \(x^n\) trong hạng đó là \(k^n/n!\). Cộng các
hạng lại là được.

Cụ thể, dùng nhân đa thức \(\Oh(c^2)\) để khai triển \(G(x)\) thành
\[
  G(x)=\sum_{k=-c}^{c} p_k e^{kx}.
\]
Khi đó xác suất là
\[
  \C{c}{m}\sum_{k=-c}^{c} p_k\left(\frac{k}{c}\right)^n,
\]
hoặc viết theo cặp \(k,-k\):
\[
  \C{c}{m}\sum_{k=1}^{c}
  \left(\frac{k}{c}\right)^n\left(p_k+(-1)^n p_{-k}\right),
\]
với hạng \(k=0\) chỉ cần xét riêng khi \(n=0\).

Độ chính xác đề yêu cầu không cao, nên các hệ số trong quá trình khai triển có
thể dùng số thực. Lũy thừa \((k/c)^n\) có thể tính bằng logarit, và khi
\(n\) lớn, chẳng hạn \(n\ge 2000\), nếu \(k\ne c\) thì \((k/c)^n\) gần như
bằng \(0\). Như vậy bài toán được giải bằng hàm sinh với độ phức tạp
\(\Oh(c^2)\).

\paragraph{Tiểu kết.}
Đây là một bài điển hình dùng hàm sinh mũ để giải bài toán có xét thứ tự. Khi
truy hồi không xử lý tốt giới hạn dữ liệu, hàm sinh mũ giải được bài này một
cách gọn gàng, một lần nữa thể hiện khả năng tối ưu hóa thuật toán của hàm
sinh. Cách dùng hàm sinh ở đây dựa trên một phép chuyển hóa nhỏ của bài toán;
năng lực chuyển hóa như vậy cũng rất đáng chú ý.

\subsection{Sweet}

\term{Ví dụ 3.}
CEOI 2004, \textit{Sweet}.

John có \(n\) lọ kẹo. Các lọ khác nhau chứa các loại kẹo khác nhau; trong cùng
một lọ thì kẹo cùng loại. Lọ thứ \(i\) có \(m_i\) viên. John quyết định ăn một
số viên kẹo; cậu muốn ăn ít nhất \(a\) viên nhưng không quá \(b\) viên. Vấn đề
là John chưa quyết định sẽ ăn bao nhiêu viên và mỗi loại ăn bao nhiêu viên.
Hỏi có bao nhiêu cách làm việc này?

Nhiệm vụ là viết chương trình đọc số kẹo trong mỗi lọ và hai số \(a,b\), xác
định số cách John có thể chọn kẹo để ăn, rồi in kết quả modulo \(2004\).
Giới hạn:
\[
  1\le n\le 10,\quad 0\le a\le b\le 10000000,\quad
  0\le m_i\le 1000000.
\]

\paragraph{Phân tích.}
Bài toán yêu cầu số cách ăn từ \(a\) đến \(b\) viên, điều này không thật dễ xử
lý trực tiếp. Ta chuyển hóa: gọi \(f(m)\) là số cách ăn nhiều nhất \(m\) viên.
Khi đó đáp án là
\[
  f(b)-f(a-1).
\]
Vì vậy chỉ cần biết cách tính \(f(m)\).

Đây là bài toán tổ hợp. Với kẹo loại \(i\), có thể ăn từ \(0\) đến \(m_i\)
viên, nên hàm sinh của nó là
\[
  G_i(x)=1+x+x^2+\cdots+x^{m_i}
        =\frac{1-x^{m_i+1}}{1-x}.
\]
Hàm sinh của toàn bộ bài toán là
\[
\begin{aligned}
  G(x)
  &=G_1(x)G_2(x)\cdots G_n(x)\\
  &=\frac{(1-x^{m_1+1})(1-x^{m_2+1})\cdots(1-x^{m_n+1})}
          {(1-x)^n}\\
  &=(1-x^{m_1+1})\cdots(1-x^{m_n+1})
    \left(1+\C{n}{n-1}x+\C{n+1}{n-1}x^2
    +\C{n+2}{n-1}x^3+\cdots\right).
\end{aligned}
\]
Trước hết cần khai triển tử số
\[
  (1-x^{m_1+1})(1-x^{m_2+1})\cdots(1-x^{m_n+1}).
\]
Mỗi ngoặc chỉ có hai hạng, nên việc khai triển \(n\) ngoặc này cần
\(\Oh(2^n)\) thời gian. Sau khi khai triển ta được một đa thức có nhiều nhất
\(2^n\) hạng khác \(0\):
\[
  \sum_{k\ge0} p_k x^k.
\]

Vì \(f(m)\) là số cách ăn từ \(0\) đến \(m\) viên, nên \(f(m)\) bằng tổng các
hệ số bậc \(0\) đến \(m\) của \(G(x)\):
\[
  f(m)=\sum_{k\ge0}p_k
  \left(1+\C{n}{n-1}+\C{n+1}{n-1}+\cdots+\C{n-1+m-k}{n-1}\right).
\]
Vấn đề còn lại là tính nhanh tổng trong ngoặc. Dùng hằng đẳng thức tổ hợp quen
thuộc
\[
  \C{r}{s}=\C{r-1}{s}+\C{r-1}{s-1},
\]
ta có
\[
\begin{aligned}
  1+\C{n}{n-1}+\C{n+1}{n-1}+\cdots+\C{n-1+m-k}{n-1}
  &=\C{n+m-k}{n}.
\end{aligned}
\]
Do đó
\[
  f(m)=\sum_{k\ge0}p_k\C{n+m-k}{n},
\]
trong đó các hạng \(k>m\) được xem là \(0\).

Cuối cùng cần tính \(\C{n+m-k}{n}\). Vì \(n\) nhỏ, có thể viết trực tiếp
\[
  \C{n+m-k}{n}
  =
  \frac{(n+m-k)(n+m-k-1)\cdots(m-k+1)}
       {n(n-1)\cdots 1}.
\]
Khi tính modulo \(2004\), có thể nhân tử số đồng thời lấy modulo
\(2004\cdot n!\), rồi chia cho \(n!\). Bước này cần \(\Oh(n)\) thời gian.

Như vậy bài toán được giải bằng hàm sinh với độ phức tạp \(\Oh(2^n n)\).

Tất nhiên, bài này cũng có thể giải bằng nguyên lý bao hàm--loại trừ. Khi đó
ta sẽ nhận được kết luận giống cách hàm sinh; phần sau của hai lời giải gần
như tương tự.

\paragraph{Tiểu kết.}
Đây là một bài đếm tổ hợp, và hàm sinh giải quyết nó rất tốt. Bài toán cũng
cho thấy hàm sinh có phạm vi ứng dụng rộng trong tổ hợp. Dù các lời giải đặc
biệt, đẹp mắt luôn hấp dẫn, phương pháp tổng quát hơn vẫn đáng quý vì nó có
thể giải nhiều bài hơn và thường có độ khó tư duy thấp hơn.

\subsection{Một bài chứng minh}

\term{Ví dụ 4.}
Dùng hàm sinh chứng minh rằng trong đồ thị đầy đủ \(n\) đỉnh có
\(n^{n-2}\) cây khung.

\paragraph{Phân tích.}
Gọi \(t_n\) là số cây khung của đồ thị đầy đủ \(n\) đỉnh, tức số cây gắn nhãn
trên \(n\) đỉnh.

Xét trường hợp tổng quát với \(n\) đỉnh. Nếu trực tiếp xét quan hệ giữa từng
cặp đỉnh, bài toán sẽ trở nên khá phức tạp. Ta dùng một thủ pháp quen thuộc:
xét riêng một đỉnh, gọi là \(p\). Sau khi bỏ đỉnh này khỏi cây, cây bị tách
thành \(m\) cây con. Nếu kích thước các cây con lần lượt là
\(k_1,k_2,\ldots,k_m\), thì số cách tạo cấu trúc bên trong mỗi cây con và nối
nó với \(p\) là
\[
  (k_1t_{k_1})(k_2t_{k_2})\cdots(k_mt_{k_m}).
\]
Kết hợp với cách chia \(n-1\) đỉnh còn lại thành các nhóm kích thước
\(k_1,\ldots,k_m\), ta được truy hồi
\[
  t_n=\sum_{m>0}\frac1{m!}
  \sum_{k_1+\cdots+k_m=n-1}
  \C{n-1}{k_1,k_2,\ldots,k_m}
  (k_1t_{k_1})\cdots(k_mt_{k_m}).
\]
Ở đây \(\C{n-1}{k_1,k_2,\ldots,k_m}\) là số cách chia \(n-1\) phần tử thành
các phần có kích thước \(k_1,\ldots,k_m\), còn \(m!\) bị chia đi vì ta không
xét thứ tự giữa các phần.

Rút gọn truy hồi:
\[
  \frac{t_n n}{n!}
  =\sum_{m>0}\frac1{m!}
   \sum_{k_1+\cdots+k_m=n-1}
   \frac{t_{k_1}k_1}{k_1!}\cdots
   \frac{t_{k_m}k_m}{k_m!}.
\]
Đặt \(g_n=nt_n\), và gọi
\[
  G(x)=\sum_{n\ge1}g_n\frac{x^n}{n!}
\]
là hàm sinh mũ của dãy \(\seq{g_n}\). Khi đó
\[
\begin{aligned}
  G(x)
  &=x\sum_{m>0}\frac1{m!}G(x)^m\\
  &=x e^{G(x)}.
\end{aligned}
\]
Đẳng thức cuối dùng khai triển
\[
  e^x=1+x+\frac{x^2}{2!}+\frac{x^3}{3!}+\cdots.
\]

Từ phương trình hàm \(G(x)=xe^{G(x)}\), dùng công thức Lagrange cho hàm ngược
ta có
\[
  [x^n]G(x)=\frac{n^{n-1}}{n!}.
\]
Vì \(G(x)\) là hàm sinh mũ, suy ra \(g_n=n^{n-1}\). Do \(g_n=nt_n\),
\[
  t_n=n^{n-2}.
\]
Trường hợp \(n=1\) cần xét riêng, vì quá trình thiết lập truy hồi ở trên giả
thiết đã bỏ một đỉnh và còn các thành phần phía sau. Mệnh đề được chứng minh.

\paragraph{Tiểu kết.}
Bài này thể hiện vai trò của hàm sinh trong chứng minh các kết luận liên quan
tới thi tin học. Hàm sinh thường được dùng để chứng minh các đẳng thức và kết
luận tổ hợp.

\subsection{Polygon}

\term{Ví dụ 5.}
IOI 2003, hoạt động thảo luận bài khó của đội tuyển quốc gia, bài 0015:
\textit{Polygon}.

Cho \(n,m\) với \(n\ge m\). Từ các đỉnh của đa giác đều \(n\) cạnh đơn vị,
chọn ra \(m\) đỉnh. Có thể tạo được bao nhiêu loại đa giác \(m\) cạnh khác
nhau? Nếu một đa giác \(m\) cạnh có thể thu được từ đa giác khác bằng phép
quay, lật hoặc tịnh tiến, thì hai đa giác được xem là cùng loại.

\paragraph{Phân tích.}
Với bài toán yêu cầu đếm số phương án khác nhau bản chất dưới phép quay và
lật, phản xạ đầu tiên là dùng định lý Pólya. Nhưng điều kiện tịnh tiến gây
khó khăn lớn cho việc xây dựng nhóm hoán vị. Ta có thể chứng minh rằng không
tồn tại hai đa giác \(m\) cạnh được tạo từ các đỉnh của đa giác đều \(n\)
cạnh mà sau tịnh tiến lại trùng nhau.

Chứng minh như sau. \(n\) đỉnh của đa giác đều cùng nằm trên một đường tròn.
Nếu có hai đa giác \(m\) cạnh trùng nhau sau tịnh tiến, thì trên đường tròn
đó sẽ có \(m\) dây cung bằng nhau và song song. Nhưng một họ đường thẳng song
song cắt một đường tròn nhiều nhất chỉ tạo ra hai dây cung bằng nhau. Do
\(m\ge3\), mâu thuẫn. Vì vậy không có hai đa giác như vậy trùng nhau sau tịnh
tiến.

Ta chỉ cần xét hai điều kiện quay và lật.

Bài toán trở thành: tô \(n\) đỉnh của đa giác đều bằng hai màu \(b,w\), trong
đó màu \(b\) được dùng đúng \(m\) lần; hỏi có bao nhiêu phương án khác nhau
bản chất dưới quay và lật. Bài toán này không thể dùng trực tiếp định lý
Pólya thường, mà cần dùng định lý Pólya dạng hàm sinh ở trên.

Ta chỉ cần xác định, với mỗi hoán vị \(a\) trong nhóm, các giá trị
\(c_1(a),\ldots,c_n(a)\).

\paragraph{1. Quay.}
Trường hợp quay khá đơn giản. Khi quay \(i\) bước, với
\(i=0,1,2,\ldots,n-1\), đặt \(d=\gcd(n,i)\). Khi đó hoán vị gồm \(d\) chu
trình, mỗi chu trình có độ dài \(n/d\), tức chỉ có
\[
  c_{n/d}(a_i)=d
\]
khác \(0\). Quy ước \(\gcd(n,0)=n\).

\paragraph{2. Lật.}
Khi lật cần phân theo tính chẵn lẻ của \(n\).

Nếu \(n\) lẻ, mỗi trục đối xứng đi qua một đỉnh và tâm đa giác. Với mỗi phép
lật \(b_i\), ta có
\[
  c_1(b_i)=1,\qquad c_2(b_i)=\frac{n-1}{2},
\]
các giá trị khác bằng \(0\).

Nếu \(n\) chẵn, có hai loại trục đối xứng:
\begin{itemize}
  \item Trục đi qua hai đỉnh đối diện: có
  \(c_1(b_i)=2\), \(c_2(b_i)=\frac n2-1\), các giá trị khác bằng \(0\).
  \item Trục đi qua trung điểm của hai cạnh đối diện: có
  \(c_2(b_i)=\frac n2\), các giá trị khác bằng \(0\).
\end{itemize}

Áp dụng định lý Pólya dạng hàm sinh, ta được:

Nếu \(n\) lẻ,
\[
  P=\frac1{2n}\left[
    \sum_{i=0}^{n-1}
      \left(b^{n/\gcd(n,i)}+w^{n/\gcd(n,i)}\right)^{\gcd(n,i)}
    +n(b+w)(b^2+w^2)^{(n-1)/2}
  \right].
\]

Nếu \(n\) chẵn,
\[
\begin{aligned}
  P=\frac1{2n}\Bigg[
    &\sum_{i=0}^{n-1}
      \left(b^{n/\gcd(n,i)}+w^{n/\gcd(n,i)}\right)^{\gcd(n,i)}\\
    &+\frac n2(b+w)^2(b^2+w^2)^{n/2-1}
     +\frac n2(b^2+w^2)^{n/2}
  \Bigg].
\end{aligned}
\]
Đáp án cần tìm là hệ số của \(b^m w^{n-m}\) trong \(P\). Chỉ cần dùng định lý
nhị thức để khai triển từng phần.

Với phần quay
\[
  \left(b^{n/d}+w^{n/d}\right)^d,\qquad d=\gcd(n,i),
\]
nếu \(n/d\) không chia hết \(m\), hệ số của \(b^mw^{n-m}\) là \(0\); ngược
lại hệ số là
\[
  \C{d}{m/(n/d)}.
\]

Với phần lật khi \(n\) lẻ,
\[
  (b+w)(b^2+w^2)^{(n-1)/2},
\]
nếu \(m\) chẵn thì hệ số là
\[
  \C{(n-1)/2}{m/2};
\]
nếu \(m\) lẻ thì hệ số là
\[
  \C{(n-1)/2}{(m-1)/2}
  =\C{(n-1)/2}{(n-m)/2}.
\]

Với phần lật qua hai đỉnh khi \(n\) chẵn,
\[
  (b+w)^2(b^2+w^2)^{n/2-1},
\]
đặt \(r=n/2-1\). Nếu \(m\) chẵn, hệ số là
\[
  \C{r}{m/2}+\C{r}{m/2-1};
\]
nếu \(m\) lẻ, hệ số là
\[
  2\C{r}{(m-1)/2}.
\]

Với phần lật qua hai cạnh khi \(n\) chẵn,
\[
  (b^2+w^2)^{n/2},
\]
nếu \(m\) chẵn thì hệ số là
\[
  \C{n/2}{m/2};
\]
nếu \(m\) lẻ thì hệ số là \(0\).

Như vậy đáp án đã được biểu diễn thành tổng của một số số tổ hợp. Ta chỉ cần
tính các số tổ hợp đó; khi \(n,m\) lớn thì dùng số nguyên lớn.

\paragraph{Tiểu kết.}
Đây là bài điển hình dùng định lý Pólya dạng hàm sinh. Nó cho thấy dạng hàm
sinh có thể tính được số phương án cụ thể hơn so với Pólya thường. Một điểm
quan trọng khác của bài là ngay từ đầu phải loại bỏ tình huống tịnh tiến; nếu
không có bước phân tích đó, bài toán rất khó bắt đầu. Điều này nhắc ta rằng
trước khi làm, cần phân tích sâu đề bài.

\section{Tổng kết}

Bài viết bắt đầu từ một số định nghĩa của hàm sinh, giới thiệu kiến thức cơ
bản và dùng ví dụ để minh họa ứng dụng của hàm sinh trong thi tin học. Đến
đây có thể thấy hàm sinh là một công cụ mạnh: trong thi tin học, nó có thể
giúp ta tối ưu hóa thuật toán và giải quyết bài toán hiệu quả. Muốn dùng tốt
hàm sinh cần có nền tảng toán học vững và năng lực phân tích tốt, để có thể
chuyển bài toán về mô hình hàm sinh. Tinh túy của tư tưởng hàm sinh là chuyển
dãy số thành hàm; tư tưởng chuyển hóa này cũng đáng được vận dụng trong các
lĩnh vực khác.

\section*{Tài liệu tham khảo}
\addcontentsline{toc}{section}{Tài liệu tham khảo}

\begin{enumerate}
  \item Wu Wenhu, Wang Jiande. \textit{Hướng dẫn Olympic Tin học: thuật toán
  và lập trình tổ hợp}. Nhà xuất bản Đại học Thanh Hoa, 2002.
  \item Liu Rujia, Huang Liang. \textit{Nghệ thuật thuật toán và thi tin học}.
  Nhà xuất bản Đại học Thanh Hoa, 2004.
  \item Ronald L. Graham, Donald E. Knuth, Oren Patashnik. \textit{Concrete
  Mathematics: A Foundation for Computer Science}, 2nd edition.
  Addison-Wesley Professional, 1994.
  \item Lu Kaicheng. \textit{Toán tổ hợp}, bản thứ ba. Nhà xuất bản Đại học
  Thanh Hoa, 2002.
  \item Zhou Yuan. Bài giảng trại mùa đông NOI, 2008.
  \item Lei Huanzhong, Jin Kai. Báo cáo lời giải hoạt động thảo luận bài khó
  của đội tuyển quốc gia IOI Trung Quốc, 2003.
\end{enumerate}

\appendix
\section{Nguyên đề của ví dụ 2}

\term{Problem F: Chocolate.}
Tệp vào: \texttt{chocolate.in}.

Năm 2100, sô-cô-la ACM sẽ là một trong những món ăn được yêu thích nhất thế
giới. Lớp vỏ đường nhiều màu như xanh lá, cam, nâu, đỏ,\ldots có lẽ là đặc
điểm hấp dẫn nhất của sô-cô-la ACM. Hiện nay người ta nói ACM chọn từ bảng
màu \(24\) màu để tô những viên kẹo ngon của họ.

Một ngày nọ, Sandy chơi một trò chơi với một gói lớn sô-cô-la ACM có năm màu
(xanh lá, cam, nâu, đỏ và vàng). Mỗi lần cậu lấy một viên từ gói và đặt lên
bàn. Nếu trên bàn có hai viên cùng màu, cậu ăn cả hai. Cậu nhận thấy một điều
khá thú vị: phần lớn thời gian trên bàn luôn có \(2\) hoặc \(3\) viên.

Bài toán đặt ra là: nếu trong gói có \(C\) màu sô-cô-la ACM, các màu phân bố
đều, sau khi lấy \(N\) viên khỏi gói thì xác suất để trên bàn có đúng \(M\)
viên là bao nhiêu?

\paragraph{Input.}
Tệp vào gồm nhiều bộ dữ liệu, mỗi bộ trên một dòng. Mỗi bộ gồm ba số nguyên
không âm \(C\), \(N\), \(M\), với \(C\le 100\), \(N,M\le 1000000\). Dữ liệu
kết thúc bởi một dòng chỉ chứa một số \(0\).

\paragraph{Output.}
Với mỗi bộ dữ liệu, in một số thực là xác suất cần tìm, làm tròn tới ba chữ
số sau dấu thập phân.

\paragraph{Ví dụ.}
\begin{verbatim}
Input
5 100 2
0

Output
0.625
\end{verbatim}

\section{Nguyên đề của ví dụ 3}

\term{Task: SWE -- Sweets.}
Ngày 1. Tệp nguồn \texttt{swe.*}. Bộ nhớ \(64\) MB. Thời gian tối đa \(0.1\)
giây.

John có \(n\) lọ kẹo. Mỗi lọ chứa một loại kẹo khác nhau; tức là kẹo trong
cùng một lọ là cùng loại, còn kẹo từ các lọ khác nhau là khác loại. Lọ thứ
\(i\) chứa \(m_i\) viên kẹo. John quyết định ăn một số viên. Cậu muốn ăn ít
nhất \(a\) viên nhưng không quá \(b\) viên. Vấn đề là cậu không quyết định
được sẽ ăn bao nhiêu viên và thuộc những loại nào. Có bao nhiêu cách để làm
việc này?

\paragraph{Nhiệm vụ.}
Viết chương trình:
\begin{itemize}
  \item đọc từ chuẩn vào số kẹo trong mỗi lọ và các số nguyên \(a,b\);
  \item xác định số cách John có thể chọn kẹo để ăn, thỏa điều kiện trên;
  \item ghi kết quả ra chuẩn ra.
\end{itemize}

\paragraph{Input.}
Dòng đầu chứa ba số nguyên \(n,a,b\), cách nhau bởi một dấu cách:
\[
  1\le n\le 10,\qquad 1\le a\le b\le 10000000.
\]
Mỗi dòng trong \(n\) dòng tiếp theo chứa một số nguyên. Dòng \(i+1\) chứa
\(m_i\), số kẹo trong lọ thứ \(i\), với \(0\le m_i\le 1000000\).

\paragraph{Output.}
Gọi \(k\) là số cách khác nhau để John chọn kẹo sẽ ăn. Dòng duy nhất của đầu
ra chứa số nguyên \(k\bmod 2004\).

\paragraph{Ví dụ.}
\begin{verbatim}
Input
2 1 3
3
5

Output
9
\end{verbatim}

John có thể chọn kẹo theo các cách:
\[
  (1,0),\ (2,0),\ (3,0),\ (0,1),\ (0,2),\ (0,3),\ (1,1),\ (1,2),\ (2,1).
\]

\end{document}
